\documentclass[10pt,a4paper]{article}
\usepackage[margin=1.9cm]{geometry}
\usepackage[T1]{fontenc}
\usepackage{lmodern}
\usepackage{amsmath}
\usepackage{xcolor}
\usepackage{booktabs}
\usepackage{array}
\usepackage{enumitem}
\usepackage{titlesec}
\usepackage{fancyvrb}
\usepackage{longtable}

\definecolor{fpgablue}{RGB}{13,71,161}
\definecolor{hostgreen}{RGB}{27,94,32}
\definecolor{accent}{RGB}{183,28,28}
\definecolor{rulegray}{RGB}{120,120,120}
\definecolor{codebg}{RGB}{245,245,245}

\usepackage[colorlinks=true]{hyperref}
\hypersetup{linkcolor=fpgablue, urlcolor=blue,
  pdftitle={HFT Project --- CDC IP Creation \& Clock Groups Runbook}}

\setlength{\parindent}{0pt}
\setlength{\parskip}{4pt}
\titleformat{\section}{\large\bfseries\color{fpgablue}}{\thesection.}{0.5em}{}
\titleformat{\subsection}{\normalsize\bfseries\color{fpgablue}}{}{0em}{}
\titlespacing*{\section}{0pt}{10pt}{4pt}
\titlespacing*{\subsection}{0pt}{6pt}{2pt}
\newcommand{\hr}{\par\textcolor{rulegray}{\rule{\linewidth}{1pt}}\par}
\newcommand{\code}[1]{\texttt{\small #1}}

% CHECKPOINT callout
\newcommand{\checkpoint}[2]{%
  \par\smallskip\noindent
  \fcolorbox{rulegray}{codebg}{%
    \begin{minipage}{\dimexpr\linewidth-2\fboxsep-2\fboxrule}
    \textcolor{hostgreen}{\textbf{CHECKPOINT #1}} --- #2
    \end{minipage}}\par\smallskip}

% WARNING callout
\newcommand{\warn}[1]{%
  \par\smallskip\noindent
  \fcolorbox{accent}{codebg}{%
    \begin{minipage}{\dimexpr\linewidth-2\fboxsep-2\fboxrule}
    #1
    \end{minipage}}\par\smallskip}

\DefineVerbatimEnvironment{cmd}{Verbatim}
  {fontsize=\small, frame=single, rulecolor=\color{rulegray},
   framesep=4pt, xleftmargin=0pt, samepage=false}

\begin{document}

{\LARGE\bfseries\color{fpgablue} CDC IP Creation \& Clock Groups}\hfill
{\small\color{rulegray} Sequential runbook \textbullet\ July 2026}\\[-1pt]
{\large From sim-verified \code{rtl/cdc} sources to five packaged IPs, integrated
across three clock domains, with a clean \code{report\_cdc}.}

\hr

You drive Vivado --- every step below is a command you run or a box you tick. Each
part ends with a \textbf{checkpoint}; do not proceed past a failed one. Companion
docs: \code{docs/ip\_packaging\_runbook.md} (original GUI walkthrough) and
\code{rtl/cdc/README\_cdc.md} (XPM production-swap table).

\begin{center}
\small
\begin{tabular}{@{}clc@{}}
\toprule
\textbf{Part} & \textbf{What you do} & \textbf{Vivado} \\
\midrule
0 & Prerequisites, part choice, and a mandatory XDC fix & no \\
1 & Package the five IPs by hand in the GUI & yes \\
2 & Register the IP repository & yes \\
3 & Verify each IP out-of-context (\code{report\_cdc} per block) & yes \\
4 & Define the three clock domains and their groups & yes \\
5 & Integrate and verify the full design & yes \\
6 & Trap reference & --- \\
A & Module port/parameter reference for IP Integrator wiring & --- \\
\bottomrule
\end{tabular}
\end{center}

\hr

\section{Prerequisites and mandatory fixes}

Do all of Part~0 before opening Vivado. Section~0.3 in particular is a real defect:
left alone it silently produces an unconstrained design that still looks green.

\subsection{0.1 Confirm the library still passes simulation}

\begin{cmd}
cd sim && make cdc-regression
\end{cmd}

\checkpoint{0.1}{ends with \code{==== CDC regression PASSED ====}, 8 tests across all
five blocks. If red, fix before packaging --- do not bake a broken block into an IP.}

\warn{\textbf{What this proves --- and does not.} \code{cdc-regression} runs under
Verilator, a \emph{zero-delay cycle simulator}. It proves protocol and data
integrity across unrelated clocks. It \textbf{cannot} model metastability or a
setup/hold violation. Green here does \emph{not} mean the crossing is silicon-safe
--- that is what Part~3 is for. Keep these two claims separate when describing this
work.}

\subsection{0.2 Pick the right part --- \code{xcvu47p} is not installed}

You will select a part in the \textbf{New Project} wizard five times (once per
block), so decide now. Use:

\begin{cmd}
xcu200-fsgd2104-2-e
\end{cmd}

\warn{Do \textbf{not} pick the F2 part \code{xcvu47p-fsvh2892-2L-e-es1}. Your local
Vivado 2025.2 install has \textbf{78 parts}, families \code{virtexuplus},
\code{zynq}, \code{azynq} only --- \code{xcvu47p} is \textbf{not among them}. It
exists only on the AWS FPGA Developer AMI build host, so it will not appear in the
wizard's part list no matter how you filter.}

\code{xcu200-fsgd2104-2-e} is the part every existing synth result in
\code{results/} was produced against, so your numbers stay comparable. In the
wizard: \textbf{Parts} tab, type \code{xcu200} in the search box, select the single
match.

To confirm from the Tcl Console at any time:

\begin{cmd}
llength [get_parts xcu200-fsgd2104-2-e]   ;# expect 1
llength [get_parts xcvu47p*]              ;# expect 0 locally
\end{cmd}

When you later package on the AWS host, choose the F2 part there instead. The CDC
logic is device-agnostic; only the packaging target changes.

\emph{Aside:} the scripted alternative \code{ip/package\_cdc\_ips.tcl} hard-codes
\code{xcvu47p} on line~11 and therefore aborts at \code{create\_project} on this
machine. You are doing this manually, so it does not affect you --- but if you ever
fall back to the script, fix that line first.

\subsection{0.3 Fix four XDC files whose constraints match zero cells}

This one is silent, which makes it the dangerous one. Several constraints select
cells with a filter like:

\begin{cmd}
[get_cells -hierarchical -filter {NAME =~ *wq_reg[0]*}]
\end{cmd}

In Tcl's \code{string match}, \textbf{\code{[0]} is a character class}, not a literal
bracket. It matches the single character \code{0}, so the pattern expects
\code{wq\_reg} followed immediately by the character \code{0}. The actual synthesized
cell is named \verb|wq_reg[0][3]| --- the next character is \verb|[|.
\textbf{The pattern matches nothing.} Verified with plain Tcl:

\begin{cmd}
*wq_reg[0]*     vs  wq_reg[0][3]        -> NO MATCH
*wq_reg\[0\]*   vs  wq_reg[0][3]        -> MATCH
*sync_reg[0]*   vs  u_fwd/sync_reg[0]   -> NO MATCH
*sync_reg\[0\]* vs  u_fwd/sync_reg[0]   -> MATCH
*sync_reg\[0\]* vs  u_fwd/sync_reg[1]   -> NO MATCH   (correctly selective)
\end{cmd}

\warn{An empty object list means \code{set\_max\_delay} / \code{set\_false\_path}
applies to \textbf{nothing}. Synthesis still succeeds, simulation stays green,
\code{report\_cdc} may still look reasonable --- and your crossing is unconstrained.
Escape the brackets.}

\textbf{The edits.} In \code{rtl/cdc/xdc/sync\_2ff\_ooc.xdc}:

\begin{cmd}
set_false_path -to [get_cells -hierarchical -filter {NAME =~ *sync_reg\[0\]*}]
\end{cmd}

In \code{rtl/cdc/xdc/pulse\_sync\_ooc.xdc} --- both \code{set\_max\_delay} targets:

\begin{cmd}
-to [get_cells -hierarchical -filter {NAME =~ *u_fwd*sync_reg\[0\]*}]
-to [get_cells -hierarchical -filter {NAME =~ *u_ack*sync_reg\[0\]*}]
\end{cmd}

In \code{rtl/cdc/xdc/async\_fifo\_ooc.xdc} and
\code{rtl/cdc/xdc/axis\_async\_fifo\_ooc.xdc}:

\begin{cmd}
-to [get_cells -hierarchical -filter {NAME =~ *wq_reg\[0\]*}]
-to [get_cells -hierarchical -filter {NAME =~ *rq_reg\[0\]*}]
\end{cmd}

\textbf{Not affected} --- these use no index and already match correctly: every
\code{set\_property ASYNC\_REG} line, and \code{handshake\_mcp\_ooc.xdc}'s
\code{*data\_hold\_reg*} $\rightarrow$ \code{*dst\_data\_reg*} \code{set\_max\_delay}.

\subsection{0.4 Decide your three domain frequencies}

The periods baked into the OOC XDCs are \emph{representative placeholders} for
out-of-context synthesis. They are overridden by the real clocks at integration
(Part~4), but pick your intended numbers now --- they drive each IP's
\code{FREQ\_HZ}.

\begin{center}
\small
\begin{tabular}{@{}lll@{}}
\toprule
\textbf{Domain} & \textbf{Role} & \textbf{Current OOC placeholder} \\
\midrule
Line / RX & market-data ingest & 156.25\,MHz (6.400\,ns) \\
Core & parser $\to$ book $\to$ strategy $\to$ risk & 250\,MHz (4.000\,ns) \\
Host / CSR & AXI-Lite, DMA, telemetry & 125\,MHz (8.000\,ns) \\
\bottomrule
\end{tabular}
\end{center}

\warn{\textbf{Reality check on Core.} Your closed result today is \textbf{200\,MHz}
(\code{tick\_to\_trade\_top} with the strategy selector, WNS $+1.815$\,ns). The XDCs
say 250\,MHz. Either retime for 250 or set the core domain to 200\,MHz and update
the periods --- but do not leave a constraint claiming a frequency you have not
closed.}

\checkpoint{0}{regression green; \code{PART} fixed; four XDCs escaped; three
frequencies chosen and written down.}

\section{Package the five IPs}

Each block is self-contained (\code{pulse\_sync} and \code{handshake\_mcp} include
their own copy of \code{sync\_2ff.sv}; \code{axis\_async\_fifo} includes
\code{async\_fifo.sv}), so \textbf{there is no ordering dependency} --- package them
in any order.

\textbf{You do this entirely by hand in the Vivado GUI} --- one project per block,
adding the source files yourself. Repeat steps 1.1--1.7 five times, substituting the
block name and its files from this table each pass:

\begin{center}
\small
\begin{tabular}{@{}clp{5.4cm}p{4.1cm}@{}}
\toprule
\textbf{Pass} & \textbf{Block (= top)} & \textbf{Design sources to add} &
\textbf{Constraints to add} \\
\midrule
1 & \code{sync\_2ff} & \code{rtl/cdc/sync\_2ff.sv} &
\code{sync\_2ff\_ooc.xdc} \\
2 & \code{pulse\_sync} & \code{rtl/cdc/pulse\_sync.sv} \newline
\code{rtl/cdc/sync\_2ff.sv} & \code{pulse\_sync\_ooc.xdc} \\
3 & \code{async\_fifo} & \code{rtl/cdc/async\_fifo.sv} &
\code{async\_fifo\_ooc.xdc} \\
4 & \code{axis\_async\_fifo} & \code{rtl/cdc/axis\_async\_fifo.sv} \newline
\code{rtl/cdc/async\_fifo.sv} & \code{axis\_async\_fifo\_ooc.xdc} \\
5 & \code{handshake\_mcp} & \code{rtl/cdc/handshake\_mcp.sv} \newline
\code{rtl/cdc/sync\_2ff.sv} & \code{handshake\_mcp\_ooc.xdc} \\
\bottomrule
\end{tabular}
\end{center}

All constraint files live in \code{rtl/cdc/xdc/}. Blocks that instantiate a
sub-module need \emph{both} files --- miss the second one and synthesis fails with
\code{ERROR: [Synth 8-439] module '<name>' not found}.

\warn{\textbf{Add the sub-module as a \emph{file}, not as a packaged IP.} When
\code{pulse\_sync} instantiates \code{sync\_2ff}, Vivado needs the RTL text of
\code{sync\_2ff} in \emph{pulse\_sync's own source list}. It does not resolve an RTL
module instantiation against your IP catalog, so the \code{sync\_2ff} IP you
packaged in pass~1 is irrelevant here --- add \code{rtl/cdc/sync\_2ff.sv} to the
project like any other source.

That means \code{sync\_2ff.sv} ends up inside three IPs and \code{async\_fifo.sv}
inside two. Deliberate: every IP is \textbf{self-contained}, so dropping any one of
them into a design pulls in everything it needs, and the five passes have no
ordering dependency at all.}

\subsection{1.1 Create the project}

\begin{enumerate}[leftmargin=1.4em, itemsep=1pt]
\item \textbf{File $\to$ Project $\to$ New\ldots} $\to$ \emph{Next}.
\item \textbf{Project name}: \code{pkg\_<block>} (e.g.\ \code{pkg\_sync\_2ff}).
  \textbf{Location}: a \code{projects/} folder that is \textbf{not} inside the IP
  location you will choose in 1.5 --- e.g.\ \code{<ip-root>/projects/}. These are
  scratch projects, not deliverables. Tick \emph{Create project subdirectory}.
  $\to$ \emph{Next}.
  \\[2pt]
  \emph{Nesting the project inside the IP folder is what produces
  \code{19-3898} (``file is not found in the project sources'') and \code{19-3656}
  (``repository path may become invalid'').} Vivado wants the two
  \textbf{adjacent}, e.g.:

\begin{cmd}
IPs/projects/pkg_pulse_sync/   <- the Vivado project
IPs/cdc_repo/pulse_sync/       <- the packaged IP (component.xml)
\end{cmd}
\item \textbf{Project type}: \emph{RTL Project}. Tick \emph{Do not specify sources
  at this time} --- you will add them deliberately in 1.2. $\to$ \emph{Next}.
\item \textbf{Default Part}: \emph{Parts} tab, search \code{xcu200}, select
  \code{xcu200-fsgd2104-2-e} (see Part~0.2). $\to$ \emph{Next} $\to$ \emph{Finish}.
\end{enumerate}

\subsection{1.2 Add the design sources by hand}

\begin{enumerate}[leftmargin=1.4em, itemsep=1pt]
\item Flow Navigator $\to$ \textbf{PROJECT MANAGER $\to$ Add Sources} (or
  \textbf{Alt+A}).
\item Select \emph{Add or create design sources} $\to$ \emph{Next}.
\item \textbf{Add Files\ldots} $\to$ browse to \code{rtl/cdc/} $\to$ select the
  \code{.sv} file(s) for this pass from the table above (Ctrl-click for two).
\item Tick \textbf{Copy sources into project}. \textbf{Do not skip this.} If left
  unticked the project references the files in place
  (\code{\$PPRDIR/../../Projects/HFT/rtl/cdc/...}) while the packager makes its own
  copies under \code{<ip-root>/src/}. The two paths do not match, so on merge the
  packager reports \code{19-3898} and \textbf{drops the file from the packaged IP}.
  A project with no \code{.srcs/} directory is the tell that this was missed.
\item \emph{Finish}.
\item In the \textbf{Sources} pane, confirm each file's type is
  \textbf{SystemVerilog}. If one came in as plain Verilog, right-click it $\to$
  \textbf{Set File Type\ldots} $\to$ \emph{SystemVerilog}. These files use
  \code{logic}, \code{always\_ff} and packed arrays --- as Verilog-2001 they will
  throw syntax errors.
\end{enumerate}

\subsection{1.3 Set the top module}

In \textbf{Sources $\to$ Design Sources}, right-click the block's module $\to$
\textbf{Set as Top}. The top module's name renders in \textbf{bold}.

\emph{If \textbf{Set as Top} does not appear, the module is already top} --- Vivado
hides the option for the module that already holds it, and with only one module in
the project it auto-designates that one. On pass~1 (\code{sync\_2ff}) that is always
the case, so there is nothing to do. Treat this step as a \emph{verification}: it
only bites on passes 2, 4 and 5, where you must confirm the \emph{wrapper} is top
and not the sub-module it instantiates. (Two other causes of a missing option:
right-clicking in the \textbf{Libraries} or \textbf{Compile Order} tab instead of
\textbf{Hierarchy}, or a file that failed to parse and landed under \emph{Non-module
Files}.)

\subsection{1.4 Add the OOC constraints}

\begin{enumerate}[leftmargin=1.4em, itemsep=1pt]
\item \textbf{Add Sources} again $\to$ \emph{Add or create constraints} $\to$
  \emph{Next}.
\item \textbf{Add Files\ldots} $\to$ \code{rtl/cdc/xdc/<block>\_ooc.xdc} $\to$ tick
  \textbf{Copy constraints files into project} $\to$ \emph{Finish}.
\item Select the XDC under \textbf{Sources $\to$ Constraints}. In the
  \textbf{Source File Properties} pane find \textbf{USED\_IN} and tick all three:
  \textbf{synthesis}, \textbf{implementation}, \code{out\_of\_context}.
\end{enumerate}

\warn{Step 3 is not optional. Without \code{out\_of\_context} in USED\_IN, the CDC
exceptions in that XDC are ignored whenever the IP is synthesized standalone --- and
standalone is exactly how a packaged OOC IP gets built. The IP will look fine and be
unconstrained.}

\subsection{1.5 Launch the packager}

\begin{enumerate}[leftmargin=1.4em, itemsep=1pt]
\item \textbf{Tools $\to$ Create and Package New IP\ldots} $\to$ \emph{Next}.
\item Choose \textbf{Package your current project} $\to$ \emph{Next}.
\item \textbf{IP location}: \code{<ip-root>/cdc\_repo/<block>} --- a
  \textbf{per-block subfolder} (create it if the browser does not). Keep all five
  under the same \code{cdc\_repo} parent: that single directory becomes your IP
  repository in Part~2. $\to$ \emph{Next}.
\item \emph{Finish}. The \textbf{Package IP} window opens with a left-hand step list.
\end{enumerate}

\warn{\textbf{Give every block its own folder.} If you point the IP location at a
bare directory (say \code{\textasciitilde/IPs}) rather than
\code{\textasciitilde/IPs/cdc\_repo/<block>}, Vivado writes \code{component.xml} to
the \emph{top} of it --- and the next block you package overwrites it. You finish
with one IP instead of five, and Part~2 has nothing to find. Check after each pass
that \code{component.xml} sits in a per-block subfolder, not at the repository root.}

\subsection{1.6 Work through the Package IP tabs}

Click each entry in the left panel in order. A tab showing a warning triangle needs
attention; one showing a tick is done.

\begin{itemize}[leftmargin=1.2em, itemsep=2pt]
\item \textbf{Identification} --- Vendor \code{hft}, Library \code{cdc}, Name
  \code{<block>}, Version \code{1.0}, Display name \code{<block>}, Description
  ``HFT CDC primitive: <block>'', Taxonomy \code{/HFT/CDC}.
\item \textbf{File Groups} --- \emph{the manifest, and the tab that matters most.}
  It declares which files actually ship inside the IP and in what role: the
  \textbf{Synthesis} group is the RTL used when the IP is synthesized \emph{plus
  your XDC}; \textbf{Simulation} is the RTL used in simulation. Confirm all three:
  \begin{enumerate}[leftmargin=1.4em, itemsep=0pt, topsep=2pt, label=(\alph*)]
  \item the block's \code{.sv} appears under \emph{Synthesis} \textbf{and}
    \emph{Simulation};
  \item the \code{\_ooc.xdc} appears under \emph{Synthesis} with file type XDC;
  \item on passes 2, 4 and 5, \textbf{both} \code{.sv} files are listed --- wrapper
    plus sub-module.
  \end{enumerate}
  If the tab offers \textbf{Merge changes from File Groups Wizard}, click it: the
  file list has gone stale, usually because you added the XDC in step~1.4 after
  opening the packager. That is nearly always why a missing XDC is missing.
  \\[2pt]
  Finally, check the \textbf{Is Include} column: it must be \textbf{unticked} on the
  XDC. \code{isIncludeFile} marks a file as a \emph{header} pulled in via
  \code{`include} --- meaningless for constraints, and it raises
  \code{19-4296}.
\item \textbf{Customization Parameters} --- \textbf{Merge changes from Customization
  Parameters Wizard}, then verify the block's parameters are listed and editable:
  \code{sync\_2ff} $\to$ \code{STAGES}, \code{INIT\_VAL};
  \code{pulse\_sync} $\to$ \code{STAGES};
  \code{async\_fifo} $\to$ \code{WIDTH}, \code{DEPTH}, \code{STAGES},
  \code{AF\_LEVEL}, \code{AE\_LEVEL};
  \code{axis\_async\_fifo} $\to$ \code{TDATA\_W}, \code{DEPTH}, \code{STAGES};
  \code{handshake\_mcp} $\to$ \code{WIDTH}, \code{STAGES}.
\item \textbf{Ports and Interfaces} --- the important tab. For each clock port,
  select it and set \textbf{FREQ\_HZ} (your Part~0.4 numbers) and
  \textbf{ASSOCIATED\_RESET}:
  \code{clk}$\to$\code{rst\_n}; \code{wr\_clk}$\to$\code{wr\_rst\_n};
  \code{rd\_clk}$\to$\code{rd\_rst\_n}; \code{s\_clk}$\to$\code{s\_rst\_n};
  \code{m\_clk}$\to$\code{m\_rst\_n}; \code{src\_clk}$\to$\code{src\_rst\_n};
  \code{dst\_clk}$\to$\code{dst\_rst\_n}.
  \\[2pt]
  On \code{axis\_async\_fifo} only: confirm \code{s\_axis\_*} and \code{m\_axis\_*}
  have been inferred as \textbf{AXI4-Stream} slave/master interfaces. If they have
  not, right-click in the interface list $\to$ \textbf{Auto Infer Interface} $\to$
  \emph{AXI4-Stream}, or add the bus interface manually and map the ports by name.
  Then set \textbf{ASSOCIATED\_BUSIF}: \code{s\_clk}$\to$\code{S\_AXIS},
  \code{m\_clk}$\to$\code{M\_AXIS}. Without this the block design will not let you
  drag-connect the streams.
\item \textbf{Review and Package} --- read the summary, then click
  \textbf{Package IP}.
\end{itemize}

\subsection{1.6a Ports and Interfaces in detail}

This tab declares what your ports \emph{mean} --- which are clocks, which are
resets, which belong to a bus --- so IP Integrator can auto-connect them. Without it
IPI sees anonymous wires.

\textbf{How to set a value:} right-click the \emph{interface} (not the raw port)
$\to$ \textbf{Edit Interface\ldots} $\to$ the \textbf{Parameters} tab. If the
parameter is not listed --- a freshly inferred clock has \emph{no} parameters at all
--- click \textbf{+} to add it by name, then set its value. POLARITY is a dropdown
(it carries a choice list) rather than free text. Plain data ports --- \code{d},
\code{q}, the AXIS payload signals --- need nothing.

Equivalently, from the Tcl Console (faster, and easy to verify afterwards):

\begin{cmd}
set core  [ipx::current_core]
set clkif [ipx::get_bus_interfaces clk -of_objects $core]

ipx::add_bus_parameter FREQ_HZ $clkif
set_property value 250000000 [ipx::get_bus_parameters FREQ_HZ -of_objects $clkif]

ipx::add_bus_parameter ASSOCIATED_RESET $clkif
set_property value rst_n [ipx::get_bus_parameters ASSOCIATED_RESET -of_objects $clkif]

ipx::update_checksums $core
ipx::save_core $core
\end{cmd}

A parameter that already exists needs no \code{add\_bus\_parameter} --- just
\code{set\_property value}.

\begin{center}
\small
\begin{tabular}{@{}llrll@{}}
\toprule
\textbf{Block} & \textbf{Clock} & \textbf{FREQ\_HZ} &
\textbf{ASSOCIATED\_RESET} & \textbf{ASSOCIATED\_BUSIF} \\
\midrule
\code{sync\_2ff}        & \code{clk}     & 250\,000\,000 & \code{rst\_n}     & --- \\
\addlinespace
\code{pulse\_sync}      & \code{src\_clk} & 125\,000\,000 & \code{src\_rst\_n} & --- \\
                        & \code{dst\_clk} & 250\,000\,000 & \code{dst\_rst\_n} & --- \\
\addlinespace
\code{async\_fifo}      & \code{wr\_clk} & 156\,250\,000 & \code{wr\_rst\_n} & --- \\
                        & \code{rd\_clk} & 250\,000\,000 & \code{rd\_rst\_n} & --- \\
\addlinespace
\code{axis\_async\_fifo} & \code{s\_clk} & 156\,250\,000 & \code{s\_rst\_n} &
\textbf{\code{S\_AXIS}} \\
                        & \code{m\_clk} & 250\,000\,000 & \code{m\_rst\_n} &
\textbf{\code{M\_AXIS}} \\
\addlinespace
\code{handshake\_mcp}   & \code{src\_clk} & 125\,000\,000 & \code{src\_rst\_n} & --- \\
                        & \code{dst\_clk} & 250\,000\,000 & \code{dst\_rst\_n} & --- \\
\bottomrule
\end{tabular}
\end{center}

The frequencies are the periods from each block's OOC XDC, so the IP metadata and
the constraints agree. If you settled on a 200\,MHz core in Part~0.4, use
\code{200000000} wherever the table says 250\,MHz.

\warn{\textbf{Set POLARITY = ACTIVE\_LOW on every reset interface.} All five blocks
use \code{always\_ff @(\ldots\ or negedge rst\_n)} --- asynchronous, active low. If
POLARITY is left ACTIVE\_HIGH, IP Integrator wires a reset that is asserted when it
should be released: the block either sits in reset forever or never resets at all.
In practice Vivado \emph{does} infer this correctly from the \code{\_n} suffix ---
on \code{sync\_2ff} the reset came through as busType \code{reset} with
\code{POLARITY = ACTIVE\_LOW} already set, so there was nothing to change. Verify
rather than assume --- the check is cheap:
\code{grep -A2 POLARITY <ip>/component.xml} should show \code{ACTIVE\_LOW}.}

\textbf{Pass 4 is the odd one.} \code{axis\_async\_fifo} is the only block with a
bus. Before setting ASSOCIATED\_BUSIF, confirm \code{s\_axis\_*} and
\code{m\_axis\_*} were inferred as \textbf{AXI4-Stream} slave/master interfaces; if
not, right-click in the interface list $\to$ \textbf{Auto Infer Interface} $\to$
\emph{AXI4-Stream}, or add the bus interface and map the ports by name. Without it
you cannot drag-connect the streams in the block design. On the other four blocks
ASSOCIATED\_BUSIF stays empty --- which is exactly why \code{19-5661} never clears
there.

\warn{\textbf{The GUI will not let you add ASSOCIATED\_BUSIF.} Its \emph{Add
Parameter} dialog offers only the parameters it considers standard for a clock
interface, and \code{ASSOCIATED\_BUSIF} is usually absent --- reported as
``ASSOCIATED\_BUSIF does not exist''. The parameter is valid; only the dialog is
limited. Add it from the Tcl Console instead:

\code{ipx::add\_bus\_parameter ASSOCIATED\_BUSIF \$clkif} then
\code{set\_property value s\_axis [ipx::get\_bus\_parameters ASSOCIATED\_BUSIF
-of\_objects \$clkif]}

Use the interface names Vivado actually created --- it derives them from the port
prefix, so they come out lowercase \code{s\_axis} / \code{m\_axis}, not the
\code{S\_AXIS} convention. Either case is accepted downstream.}

\subsection{1.6b Warnings you will see --- and which ones matter}

The \textbf{Messages} pane fills with IP Packager warnings on every pass. Most are
noise. Note that Vivado groups them by \emph{wizard}, and the badge on a left-panel
tab counts that wizard's messages --- which is not always the tab where you fix the
problem. Triage:

\begin{center}
\small
\begin{tabular}{@{}p{4.3cm}p{5.0cm}p{5.2cm}@{}}
\toprule
\textbf{Message} & \textbf{Means} & \textbf{Do} \\
\midrule
\code{19-11770} Clock interface \code{clk} has no FREQ\_HZ parameter &
the clock carries no declared frequency &
\textbf{FIX} --- Ports and Interfaces $\to$ select the clock $\to$ set
\textbf{FREQ\_HZ} (and \textbf{ASSOCIATED\_RESET}) \\
\addlinespace
\code{19-11888} IP description is not meaningful &
Description still defaults to the module name. \emph{Reported under the
Compatibility Wizard, so the badge lands on the Compatibility tab} &
\textbf{FIX} (cosmetic) --- but on the \textbf{Identification} tab, not
Compatibility: write a real description \\
\addlinespace
\code{19-5661} Bus Interface \code{clk} does not have any bus interfaces
associated with it &
\code{ASSOCIATED\_BUSIF} is empty &
\textbf{IGNORE} on \code{sync\_2ff}, \code{pulse\_sync}, \code{async\_fifo},
\code{handshake\_mcp} --- they have no AXI interface. \textbf{Only}
\code{axis\_async\_fifo} must set it \\
\addlinespace
\code{19-4296} Unsupported include file should not have \code{isIncludeFile} set &
the \emph{Is Include} box is ticked on the XDC; that flag marks a
\code{`include} header, not constraints &
\textbf{FIX} --- File Groups $\to$ untick \emph{Is Include} on the XDC, re-package \\
\addlinespace
\code{19-5101} Packaging a component with a SystemVerilog top file is not fully
supported &
the packager's SV support at top level is partial &
\textbf{IGNORE} here --- these blocks expose only \code{logic} scalars and packed
vectors, which map cleanly. Only if ports or parameters come out wrong in IPI, add a
thin Verilog-2001 wrapper for packaging \\
\addlinespace
\code{19-3656} repository path may become invalid if the project moves &
the packaging project is nested \emph{inside} the IP folder &
\textbf{IGNORE} if it already works; for later passes put the project
\emph{adjacent} to the IP (Part 1.1) \\
\addlinespace
\code{19-2187} The Product Guide file is missing &
a documentation-completeness check --- commercial IPs ship a datasheet PDF the way
Xilinx IPs ship their PG documents &
\textbf{IGNORE} --- nothing to attach on a hand-rolled block; appears on all five
passes \\
\bottomrule
\end{tabular}
\end{center}

\warn{\textbf{Do not expect a clean sheet.} \code{19-5661} and \code{19-2187} never
clear --- they persist on all five passes and that is correct behaviour, not a
defect. They are \emph{warnings}, not errors: \textbf{Package IP} stays clickable
throughout. Only \code{19-11770} (FREQ\_HZ) and \code{19-11888} (description) call
for action. Note too that these are \emph{metadata} problems --- rearranging folders
or re-creating the project will not touch them; you clear them by filling fields in
the wizard.}

\warn{\textbf{Check the Identification defaults.} Out of the box the component is
named something like \code{Twinwinglabs:user:sync\_2ff:1.0} --- your Vivado install's
default vendor --- or \code{user.org:user:sync\_2ff:1.0} on a stock install. Set
\textbf{Vendor} \code{hft}, \textbf{Library} \code{cdc} and \textbf{Taxonomy}
\code{/HFT/CDC} on every pass, or the five blocks will not group together under
\emph{HFT/CDC} in the catalog in Part~2. To check what you actually shipped, grep
the packaged component: \code{grep spirit:vendor <ip>/component.xml}, and confirm
\code{FREQ\_HZ} appears at all --- zero hits means 1.6's clock settings never took.}

\textbf{On the Compatibility tab:} leave \emph{Package for Vitis} unchecked --- that
is for Vitis acceleration kernels with an \code{ap\_ctrl\_hs} control protocol, not
what these are. Do \textbf{not} tick \emph{Ignore Freq\_Hz}; you want FREQ\_HZ.
Confirm the \textbf{Family} list includes \code{virtexuplus}: \code{xcvu47p} is a
Virtex UltraScale+ device, so even though that part is not installed locally the
family is covered and these IPs stay usable when you later build on the AWS host.

\subsection{1.7 Close and repeat}

Vivado offers to close the project after packaging --- accept. Return to 1.1 for the
next block. Because each block carries its own copy of any sub-module, \textbf{there
is no ordering dependency}: package them in whatever order you like.

\emph{Note on OOC synthesis:} the Out-Of-Context choice is \textbf{not} made in the
packaging wizard. It is chosen later, when you instantiate the IP and run
\textbf{Generate Output Products} --- set \textbf{Synthesis Options} to
\emph{Out of context per IP} so each IP synthesizes and caches independently.

\checkpoint{1}{five packaged IPs exist, one folder each, every one containing a
\code{component.xml} --- \emph{and each one passes the checker}:

\code{python3 ip/check\_packaged\_ip.py <ip-root>/*}

It reads the IP-XACT only (it never launches Vivado) and exits non-zero if any
\textbf{FAIL} remains. Trust it over the wizard's badges: it verifies what actually
got written to disk.}

\textbf{Do not skip the checker.} The GUI reports success while leaving real
problems behind --- observed on the first live pass, \emph{after} the packager said
it was done: \code{isIncludeFile} had migrated from the XDC onto \emph{both} copies
of \code{sync\_2ff.sv} (so the module would be treated as an include header rather
than compiled), and the clock interface had picked up a stray \code{POLARITY} while
still missing \code{FREQ\_HZ} and \code{ASSOCIATED\_RESET}. None of that is visible
without reading \code{component.xml}.

\warn{\textbf{When you untick \emph{Is Include}, untick it everywhere.} The flag is
per-file and per-file-group: the same \code{.sv} appears in both the Synthesis and
Simulation groups and each carries its own checkbox. Clearing it on one row --- or
clearing it on the XDC and accidentally setting it on the sources --- leaves you
exactly where you started. Re-run the checker after every re-package.}

\section{Register the IP repository}

In the project where you will assemble the design:

\begin{cmd}
set_property ip_repo_paths [file normalize ip/cdc_repo] [current_project]
update_ip_catalog
\end{cmd}

\textbf{You need a project open first.} An IP repository is a project setting, so
there is no such menu on the Welcome screen. Since Part~4 needs an integration
project anyway, create it now: \textbf{File $\to$ Project $\to$ New\ldots}, name
\code{cdc\_integration}, RTL Project, \emph{Do not specify sources at this time},
part \code{xcu200-fsgd2104-2-e}. The repository registration then carries straight
through to the block design.

\textbf{GUI route (use this one):} Flow Navigator $\to$ \textbf{PROJECT MANAGER}
$\to$ \textbf{Settings} $\to$ \textbf{IP $\to$ Repository} $\to$ the \textbf{+}
button $\to$ browse to your IP root $\to$ \emph{Select}. Vivado reports how many
IPs it found --- \textbf{it should say 5}. Click \emph{OK}, then \textbf{Refresh IP
Catalog} if prompted. (Equivalent: \textbf{Window $\to$ IP Catalog}, right-click
$\to$ \textbf{Add Repository\ldots})

\warn{\textbf{Do not use \emph{Manage IP}.} It looks like the right feature and is
not: Manage IP is a flow for generating IP customisations out-of-context, and it
demands a directory containing a \code{manage\_ip\_project}. Pointing it at your IP
root gives ``The directory you selected is not a valid Manage IP location''. Adding
a \emph{repository} is the Settings route above.}

If it finds fewer than five, one of the packaging passes did not complete; check
which \code{ip/cdc\_repo/<block>/} folder is missing its \code{component.xml} and
redo that pass.

\checkpoint{2}{all five blocks appear in \textbf{Window $\to$ IP Catalog} under
\textbf{HFT / CDC}.}

\section{Verify each IP out-of-context}

\textbf{Do this before integrating.} Per-block OOC checks prove the primitives are
sound; integration checks (Part~5) prove your \emph{wiring} is sound. Debugging both
at once is much harder. This is the verification layer that has never been run on
this library.

\subsection{3.1 Run the checks --- GUI route}

\warn{\textbf{Synthesis only --- do not run implementation.} \code{report\_cdc},
\code{report\_clock\_interaction} and \code{check\_timing} all work on the
\emph{synthesized} netlist. CDC is a structural property of how signals cross
domains, fully determined after synthesis; place \& route adds nothing to the
analysis. Packaging itself (Part~1) needs neither --- the packager reads the RTL and
metadata directly.

\textbf{And ignore per-block area and power.} \code{sync\_2ff} is two flip-flops;
its utilisation is a couple of registers and its power is rounding error. A
synthesis-only power report is vectorless and low-confidence anyway (real numbers
need post-route plus a SAIF). The metrics that mean something come from the
\emph{integrated} design in Part~5, measured against the existing baseline ---
WNS $+1.815$\,ns / 7570 LUTs at 200\,MHz for the selector top, and WNS $+0.066$\,ns
from the real \code{xcvu47p} P\&R. The question worth answering is what the CDC
restructuring costs against those, not what a synchroniser weighs alone.}

\textbf{Reuse the \code{pkg\_<block>} projects you built in Part~1.} They already
contain the right sources, the OOC XDC and the correct part, so there is nothing to
set up. For each block:

\begin{enumerate}[leftmargin=1.4em, itemsep=1pt]
\item \textbf{File $\to$ Project $\to$ Open Recent} $\to$ \code{pkg\_<block>}.
\item Flow Navigator $\to$ \textbf{SYNTHESIS $\to$ Run Synthesis}. Accept the
  defaults; these blocks synthesize in seconds.
\item When it finishes choose \textbf{Open Synthesized Design}.
\item \textbf{Reports $\to$ Timing $\to$ Report CDC\ldots} $\to$ tick
  \emph{Details} $\to$ \emph{OK}. Results open in a docked pane; use \emph{Export}
  to save to \code{results/cdc/<block>\_cdc.rpt}.
\item \textbf{Reports $\to$ Timing $\to$ Report Clock Interaction\ldots} $\to$
  \emph{OK}.
\item \textbf{Reports $\to$ Timing $\to$ Check Timing\ldots} $\to$ \emph{OK}.
\end{enumerate}

\emph{Equivalent Tcl}, if you prefer typing into the \textbf{Tcl Console} at the
bottom of the same GUI window (this is where the 3.2 checks must go anyway):

\begin{cmd}
report_cdc               -details -file results/cdc/async_fifo_cdc.rpt
report_clock_interaction -file results/cdc/async_fifo_clkint.rpt
check_timing   -verbose  -file results/cdc/async_fifo_checktiming.rpt
\end{cmd}

For reference, the sources each block needs (top module = the first file) --- the
same set you added by hand in Part~1.2:

\begin{center}
\small
\begin{tabular}{@{}ll@{}}
\toprule
\textbf{Block} & \textbf{Sources} \\
\midrule
\code{sync\_2ff}        & \code{sync\_2ff.sv} \\
\code{pulse\_sync}      & \code{pulse\_sync.sv} \quad \code{sync\_2ff.sv} \\
\code{async\_fifo}      & \code{async\_fifo.sv} \\
\code{axis\_async\_fifo} & \code{axis\_async\_fifo.sv} \quad \code{async\_fifo.sv} \\
\code{handshake\_mcp}   & \code{handshake\_mcp.sv} \quad \code{sync\_2ff.sv} \\
\bottomrule
\end{tabular}
\end{center}

\subsection{3.2 Confirm the constraints actually bound to cells}

This is the Part~0.3 bug class. \textbf{After \code{synth\_design}, before trusting
any report:}

\begin{cmd}
# how many flops actually carry ASYNC_REG?
llength [get_cells -hier -filter {ASYNC_REG}]

# do the escaped patterns now select real cells? (expect non-zero)
llength [get_cells -hierarchical -filter {NAME =~ *wq_reg\[0\]*}]
llength [get_cells -hierarchical -filter {NAME =~ *rq_reg\[0\]*}]

# and confirm the exceptions landed
llength [get_timing_exceptions]
\end{cmd}

Expected non-zero ASYNC\_REG counts: \code{sync\_2ff} = STAGES (2);
\code{pulse\_sync} and \code{handshake\_mcp} = 2 chains $\times$ STAGES (4);
\code{async\_fifo} / \code{axis\_async\_fifo} = $2 \times$ STAGES $\times$ (AW$+1$)
bits.

\warn{\textbf{If any count is 0, stop.} The constraint is not applied and every
report below is meaningless.}

\subsection{3.3 Acceptance criteria}

\textbf{\code{report\_cdc}} --- every crossing recognized and \textbf{Safe}; zero
Unsafe, zero Unknown.

\begin{center}
\small
\begin{tabular}{@{}p{3.1cm}p{11.6cm}@{}}
\toprule
\textbf{Block} & \textbf{Expect} \\
\midrule
\code{sync\_2ff} & \code{d} is an async input into a 2-FF ASYNC\_REG chain,
false-pathed. Single clock, so no clock-pair crossing. \\
\code{pulse\_sync} & Both toggle chains (\code{u\_fwd}, \code{u\_ack}) safe, 1-bit each. \\
\code{async\_fifo} & Gray-coded pointer crossings classified safe (multi-bit but
Gray, one bit changes per step). \\
\code{axis\_async\_fifo} & Same as \code{async\_fifo} --- a thin
\code{\{tlast,tdata\}} wrapper. \\
\code{handshake\_mcp} & 1-bit req/ack toggles safe; the wide \code{data\_hold}
$\to$ \code{dst\_data} bus recognized as an MCP crossing covered by
\code{set\_max\_delay -datapath\_only}. \\
\bottomrule
\end{tabular}
\end{center}

Read \code{handshake\_mcp} and the FIFOs carefully rather than just checking for
green. If the wide bus is reported as an \textbf{unsafe multi-bit crossing}, the
\code{set\_max\_delay} is not matching cell names (Part~0.3 again). The FIFO memory
may also be flagged as a multi-bit crossing --- that one is expected and justified
(reads are gated by the synchronized pointer), and it is exactly the case you should
be able to \emph{explain} rather than waive.

\textbf{\code{report\_clock\_interaction}} --- every async clock pair must read
\textbf{``Asynchronous Groups''}. If a pair reads \textbf{``Partial False Path''},
that is the classic silent bug: some paths across the boundary are constrained and
some are not. Treat it as a failure, not a warning.

\textbf{\code{check\_timing}} --- no unconstrained endpoints, no missing
\code{create\_clock}, no \code{no\_clock} warnings.

\checkpoint{3}{for all five blocks --- non-zero ASYNC\_REG counts, zero
Unsafe/Unknown in \code{report\_cdc}, no ``Partial False Path'', clean
\code{check\_timing}. Save the reports to \code{results/cdc/}: they are the evidence
that the library is sound.}

\section{Clock domains and clock groups}

\subsection{4.1 The three domains and where each crossing sits}

\begin{cmd}
   LINE domain              CORE domain                   HOST domain
  (RX / ingest)      (parse -> book -> strategy -> risk) (AXI-Lite / DMA)

  ITCH bytes --[axis_async_fifo]--> itch_parser ... risk_check
                  s_clk=line                |
                  m_clk=core                +--[axis_async_fifo]--> egress/DMA
                                                 s_clk=core, m_clk=host

  config (strat_sel, thresholds) <--[handshake_mcp]-- host CSR writes
                                                      src=host, dst=core
  halt / enable bit              <--[sync_2ff]------- host
  GO / latency-clear strobe      <--[pulse_sync]----- host
  latency snapshots              --[handshake_mcp]--> host telemetry
                                    src=core, dst=host
\end{cmd}

\warn{\textbf{Rule:} no signal crosses a boundary except through one of these five
blocks. A signal that straddles the boundary in raw logic is what turns
\code{report\_cdc} \emph{Unknown}.}

\subsection{4.2 Write the top-level clock constraints}

At integration the IP-level OOC \code{create\_clock} commands are replaced by the
real top-level clocks (shell clock recipe, an MMCM, or board pins). Define all
three, then declare all three mutually asynchronous:

\begin{cmd}
# --- top-level clock definitions (adjust to your actual sources) ------------
create_clock -name clk_line -period 6.400 [get_ports clk_line]   ;# 156.25 MHz
create_clock -name clk_core -period 4.000 [get_ports clk_core]   ;# 250 MHz (see 0.4)
create_clock -name clk_host -period 8.000 [get_ports clk_host]   ;# 125 MHz

# --- the three domains are mutually unrelated ------------------------------
set_clock_groups -asynchronous \
  -group [get_clocks clk_line] \
  -group [get_clocks clk_core] \
  -group [get_clocks clk_host]
\end{cmd}

One \code{set\_clock\_groups} with three \code{-group} options declares all three
pairs async --- do not write three separate two-group commands, which is easy to get
subtly wrong. If the clocks come from an MMCM rather than separate oscillators, use
the generated clock names (\code{get\_clocks -of\_objects [get\_pins ...]}) instead
of port names.

\warn{\textbf{If you later move to an Arty A7-100T:} that board has a
\emph{single} 100\,MHz oscillator, so MMCM-derived clocks are phase-\emph{related}
and Vivado sees them as synchronous. You can still constrain them async and exercise
the logic, but it is only a partial async-CDC demonstration --- two truly independent
oscillators are what makes this real. Also, Artix-7 device support is not installed
in your Vivado today (\code{get\_parts xc7a100t*} returns 0); you would add it via
the installer.}

\subsection{4.3 Assemble in IP Integrator}

\textbf{Create the ports first}, before adding any IP --- right-click canvas $\to$
\textbf{Create Port\ldots} (Ctrl+K). Setting Type=Clock with a frequency is what
makes IPI propagate it into each IP's FREQ\_HZ.

\begin{center}
\small
\begin{tabular}{@{}llll@{}}
\toprule
\textbf{Port} & \textbf{Direction} & \textbf{Type} & \textbf{Setting} \\
\midrule
\code{clk\_line}  & Input & Clock & 156.25\,MHz \\
\code{clk\_core}  & Input & Clock & 250\,MHz \\
\code{clk\_host}  & Input & Clock & 125\,MHz \\
\code{rst\_line\_n} & Input & Reset & ACTIVE\_LOW \\
\code{rst\_core\_n} & Input & Reset & ACTIVE\_LOW \\
\code{rst\_host\_n} & Input & Reset & ACTIVE\_LOW \\
\bottomrule
\end{tabular}
\end{center}

Then \textbf{Add IP} (right-click canvas): \code{axis\_async\_fifo} $\times 2$
(ingress + egress), \code{handshake\_mcp}, \code{sync\_2ff}, \code{pulse\_sync}.
Double-click each to set parameters --- ingress \code{TDATA\_W}=8, egress
\code{TDATA\_W}=9, \code{handshake\_mcp} \code{WIDTH}=32, \code{STAGES}=2 throughout.

\textbf{To draw a connection:} hover a pin until the cursor becomes a pencil, then
drag to the target; valid targets highlight green. Work \emph{from each port
outward} rather than block by block --- far fewer canvas jumps. Eighteen
connections in total:

\begin{center}
\small
\begin{tabular}{@{}lp{10.2cm}@{}}
\toprule
\textbf{From port} & \textbf{To pins} \\
\midrule
\code{clk\_line}   & \code{axis0/s\_clk} \\
\code{clk\_core}   & \code{axis0/m\_clk}, \code{axis1/s\_clk},
                     \code{handshake/dst\_clk}, \code{sync\_2ff/clk},
                     \code{pulse\_sync/dst\_clk} \\
\code{clk\_host}   & \code{axis1/m\_clk}, \code{handshake/src\_clk},
                     \code{pulse\_sync/src\_clk} \\
\addlinespace
\code{rst\_line\_n} & \code{axis0/s\_rst\_n} \\
\code{rst\_core\_n} & \code{axis0/m\_rst\_n}, \code{axis1/s\_rst\_n},
                      \code{handshake/dst\_rst\_n}, \code{sync\_2ff/rst\_n},
                      \code{pulse\_sync/dst\_rst\_n} \\
\code{rst\_host\_n} & \code{axis1/m\_rst\_n}, \code{handshake/src\_rst\_n},
                      \code{pulse\_sync/src\_rst\_n} \\
\bottomrule
\end{tabular}
\end{center}

\warn{\textbf{Each reset follows its clock exactly.} \code{clk\_core} and
\code{rst\_core\_n} go to the same five pins; likewise for the other two domains. A
reset landing on a pin whose clock came from a different domain is the mistake to
look for --- it is silent in the block design and only shows up as odd behaviour
later. Never fan one reset across all three domains.}

\textbf{Finally, export the data ports:} select each remaining pin $\to$
\textbf{Ctrl+T} (Make External). For AXIS, select the \emph{interface} pin (the
bundled connector) and one Ctrl+T exports the whole bus: \code{S\_AXIS} and
\code{M\_AXIS} on both FIFOs. Then individually: \code{handshake\_mcp}'s
\code{src\_valid}/\code{src\_data}/\code{src\_ready}/\code{dst\_valid}/\code{dst\_data};
\code{sync\_2ff}'s \code{d}/\code{q}; \code{pulse\_sync}'s
\code{src\_pulse}/\code{src\_busy}/\code{dst\_pulse}. Nothing may be left dangling
--- Validate fails on unconnected pins and names the offender.

\checkpoint{4}{block design validates (\textbf{Tools $\to$ Validate Design}) with no
critical warnings; every cross-domain connection passes through a CDC IP.}

\section{Verify the integrated design}

After synthesizing the assembled design:

\begin{cmd}
report_cdc               -details -file results/cdc/top_cdc.rpt
report_clock_interaction -file results/cdc/top_clkint.rpt
report_timing_summary    -file results/cdc/top_timing.rpt
check_timing   -verbose  -file results/cdc/top_checktiming.rpt
\end{cmd}

Acceptance gates:

\begin{enumerate}[leftmargin=1.4em, itemsep=2pt]
\item \textbf{\code{report\_cdc}}: zero Unsafe, zero Unknown. An \emph{Unknown}
  almost always means a signal bypassed the CDC IPs --- fix it by routing through
  the right primitive, \textbf{not} by adding a waiver.
\item \textbf{\code{report\_clock\_interaction}}: all three pairs show
  ``Asynchronous Groups''; no ``Partial False Path''.
\item \textbf{\code{report\_timing\_summary}}: each clock meets its own period
  independently; WNS positive per domain.
\item \textbf{\code{check\_timing}}: no unconstrained endpoints.
\end{enumerate}

\checkpoint{5}{all four gates pass. At that point you have a genuine three-domain
design with \emph{static CDC sign-off} --- a materially stronger claim than ``the
simulation passed''.}

\section{Trap reference}

\begin{center}
\small
\begin{tabular}{@{}cp{5.1cm}p{4.3cm}p{4.0cm}@{}}
\toprule
\textbf{\#} & \textbf{Trap} & \textbf{Symptom} & \textbf{Fix} \\
\midrule
1 & Reaching for the F2 part \code{xcvu47p} & it is not in the New Project part
list at all & Part 0.2 --- use \code{xcu200-fsgd2104-2-e} \\
2 & \code{[0]} in a \code{NAME =\textasciitilde} filter is a character class &
constraint silently applies to nothing & Part 0.3 --- escape as
\verb|\[0\]| \\
3 & XDC missing \code{USED\_IN out\_of\_context} & exceptions ignored during OOC
synth & set the file property in Part 1.3 \\
4 & \code{set\_property ASYNC\_REG} matched zero cells & synchronizers unprotected,
sim still green & count with \code{get\_cells -hier -filter \{ASYNC\_REG\}} (3.2) \\
5 & ``Partial False Path'' in clock interaction & some boundary paths unconstrained
& one \code{set\_clock\_groups} covering all pairs (4.2) \\
6 & Constraint claims 250\,MHz, design closes 200 & timing signed off against
fiction & Part 0.4 --- set the real number \\
7 & Treating Verilator green as CDC sign-off & metastability never modelled &
Part 3 is the real check \\
8 & One reset fanned across all domains & reset removal violates in two domains &
per-domain synchronized reset (4.3) \\
9 & Only the wrapper \code{.sv} added on passes 2, 4, 5 & elaboration fails,
sub-module not found & Part 1.2 --- add \emph{both} files \\
10 & \code{.sv} added as Verilog-2001, not SystemVerilog & syntax errors on
\code{logic} / \code{always\_ff} & Part 1.2 --- Set File Type $\to$ SystemVerilog \\
11 & \emph{Copy sources into project} left unticked & \code{19-3898}: file ``not
found in the project sources'', XDC silently dropped from the IP & Part 1.2 --- tick
it; no \code{.srcs/} dir is the tell \\
12 & IP location is a bare folder, not a per-block subfolder & each pass overwrites
the last \code{component.xml}; one IP instead of five & Part 1.5 --- use
\code{cdc\_repo/<block>} \\
13 & Spaces in the project path (e.g.\ \code{IP for HFT}) & Vivado copes, but Tcl
paths, simulator scripts and shell tooling break later & Part 1.1 --- use a
space-free path \\
14 & One project for all five blocks & \emph{Package your current project} packages
one IP from one top module & Part 1.1 --- one project per block \\
\bottomrule
\end{tabular}
\end{center}

\section*{Appendix A --- Module reference for IPI wiring}
\addcontentsline{toc}{section}{Appendix A}

\subsection*{\code{sync\_2ff} --- single-bit level crossing}
Params \code{STAGES=2}, \code{INIT\_VAL=0}. \\
Ports \code{clk}, \code{rst\_n}, \code{d} (async in), \code{q} (synchronized out).

\subsection*{\code{pulse\_sync} --- one-cycle strobe crossing (toggle + closed-loop ack)}
Params \code{STAGES=2}. \\
Src \code{src\_clk}, \code{src\_rst\_n}, \code{src\_pulse}, \code{src\_busy}. \\
Dst \code{dst\_clk}, \code{dst\_rst\_n}, \code{dst\_pulse}. \\
\emph{\code{src\_busy} is high while a pulse is in flight --- gate new pulses on it
or they are dropped.}

\subsection*{\code{async\_fifo} --- multi-bit stream, Gray-pointer dual-clock, FWFT}
Params \code{WIDTH=8}, \code{DEPTH=16} (power of two), \code{STAGES=2},
\code{AF\_LEVEL=DEPTH-2}, \code{AE\_LEVEL=1}. \\
Write \code{wr\_clk}, \code{wr\_rst\_n}, \code{wr\_en}, \code{wr\_data},
\code{full}, \code{almost\_full}. \\
Read \code{rd\_clk}, \code{rd\_rst\_n}, \code{rd\_en}, \code{rd\_data},
\code{empty}, \code{almost\_empty}.

\subsection*{\code{axis\_async\_fifo} --- AXI-Stream boundary crossing}
Params \code{TDATA\_W=8}, \code{DEPTH=16}, \code{STAGES=2}. \\
Slave \code{s\_clk}, \code{s\_rst\_n}, \code{s\_axis\_tvalid}, \code{s\_axis\_tready},
\code{s\_axis\_tdata}, \code{s\_axis\_tlast}. \\
Master \code{m\_clk}, \code{m\_rst\_n}, \code{m\_axis\_tvalid}, \code{m\_axis\_tready},
\code{m\_axis\_tdata}, \code{m\_axis\_tlast}. \\
\emph{Carries \code{\{tlast, tdata\}} through one \code{async\_fifo} of width
\code{TDATA\_W+1}.}

\subsection*{\code{handshake\_mcp} --- atomic multi-bit config (req/ack, data held stable)}
Params \code{WIDTH=32}, \code{STAGES=2}. \\
Src \code{src\_clk}, \code{src\_rst\_n}, \code{src\_valid}, \code{src\_data},
\code{src\_ready}. \\
Dst \code{dst\_clk}, \code{dst\_rst\_n}, \code{dst\_valid} (1-cycle pulse),
\code{dst\_data} (held stable). \\
\emph{Load only when \code{src\_ready} is high. Only the toggles cross; the data bus
is an MCP relaxed with \code{set\_max\_delay -datapath\_only}, never a synchronized
bus.}

\end{document}
